\subsection{Analisis Regulasi}
\label{subsec:analisis-regulasi}
Analisis regulasi pada penelitian ini mengacu pada pembahasan mengenai kepemilikan dan hak akses RME pada subbab \ref{subsec:kepemilikan-hak-akses-rme}. Permenkes 24/2022 menegaskan bahwa dokumen Rekam Medis Elektronik dikelola oleh fasyankes, sedangkan isi rekam medis merupakan milik pasien. Konsekuensinya, rancangan sistem harus tetap menempatkan pasien sebagai pemilik otoritas utama atas penyampaian isi rekam medis, sekaligus menyediakan mekanisme yang memungkinkan tenaga medis melaksanakan tugasnya secara sah dalam proses pelayanan.

Regulasi yang sama juga menunjukkan bahwa akses terhadap RME tidak bersifat seragam untuk semua pihak. Tenaga medis dapat melakukan aksi yang berbeda terhadap data, seperti menambahkan, memperbaiki, atau membaca isi rekam medis, sedangkan pihak lain di lingkungan fasyankes memiliki kebutuhan akses yang lebih terbatas. Ditambah lagi, UU PDP mengklasifikasikan data kesehatan sebagai Data Pribadi Spesifik yang wajib dilindungi dari pemrosesan dan akses yang tidak sah. Implikasi teknisnya adalah sistem tidak cukup hanya membedakan antara ``boleh akses'' dan ``tidak boleh akses'', tetapi perlu mendukung pembatasan yang lebih rinci berdasarkan jenis data dan tindakan yang diizinkan.

Dari sudut pandang kepatuhan, analisis regulasi tersebut menghasilkan dua kebutuhan utama.
\begin{enumerate}
    \item Sistem harus mampu menegakkan prinsip \textit{need-to-know} atau \textit{least privilege}, sehingga setiap aktor hanya memperoleh akses pada bagian data yang benar-benar relevan dengan tugasnya.
    \item Mekanisme pemberian akses perlu tetap operasional dalam alur klinis yang kolaboratif tanpa menghilangkan kendali awal pasien atas otorisasi.
\end{enumerate}
Dengan demikian, regulasi justru menguatkan kebutuhan akan kontrol akses granular, segmentasi data \textit{off-chain}, dan pendelegasian berantai yang terbatas, yang selaras dengan rumusan masalah dan tujuan Tugas Akhir ini.
